\documentclass[11pt,a4paper]{article}
\usepackage[utf8]{inputenc}
\usepackage[vietnamese]{babel}
\usepackage{amsmath,amssymb}
\usepackage{geometry}
\geometry{a4paper, total={180mm,260mm}, left=15mm, right=15mm, top=18mm, bottom=18mm}

\usepackage{tikz}
\usetikzlibrary{shapes.geometric, arrows.meta, positioning, calc}
\tikzset{
  arr/.style={-{Stealth[scale=0.8]}},
  every node/.append style={align=center}
}
\usepackage{tcolorbox}
\tcbuselibrary{skins,breakable}
\usepackage{fontawesome}
\usepackage{fancyhdr}
\usepackage{enumitem}

\pagestyle{fancy}
\fancyhf{}
\renewcommand{\headrulewidth}{0.5pt}
\renewcommand{\footrulewidth}{0.5pt}
\fancyhead[L]{\small\bfseries\color{violet!85!black}\faBook\ SỔ TAY KIẾN THỨC TOÁN 2}
\fancyhead[R]{\small\bfseries\color{violet!85!black}Lớp Toán Thầy Lĩnh -- loptoanthaylinh.vn}
\fancyfoot[C]{\small\color{gray}-- Sổ tay kiến thức: Bài 23. Phép trừ có nhớ số có 2 chữ số cho số có 2 chữ số --}

\setlength{\parindent}{0pt}
\setlength{\parskip}{0pt}

\begin{document}

% --- BANNER BÀI HỌC ---
\begin{tcolorbox}[
  enhanced,
  colback=violet!8!white,
  colframe=violet!70!black,
  arc=3mm, boxrule=1.2pt,
  left=3mm, right=3mm, top=2mm, bottom=2mm,
  halign=center
]
  {\LARGE\bfseries\color{violet!90!black}\faBookmark\ BÀI 23: PHÉP TRỪ (CÓ NHỚ) SỐ CÓ HAI CHỮ SỐ CHO SỐ CÓ HAI CHỮ SỐ}\\[1mm]
  {\itshape\footnotesize\color{black!75} Kĩ thuật trừ hai số có 2 chữ số có nhớ dạng 64 - 28 trên 1 trang A4}
\end{tcolorbox}

\vspace{-2mm}

\noindent
\begin{minipage}[t]{0.49\textwidth}\vspace{0pt}
\begin{tcolorbox}[
  enhanced,
  colback=violet!4!white,
  colframe=violet!75!black,
  arc=2mm, boxrule=0.8pt,
  left=2mm, right=2mm, top=1.8mm, bottom=1.8mm,
  title={\footnotesize\bfseries\color{white}\faKey\ 1. Quy tắc trừ 2 chữ số có nhớ}
]
{\footnotesize
\begin{itemize}\setlength{\itemsep}{1pt}\setlength{\leftskip}{-2.5mm}
  \item \textbf{Bước 1: Đặt tính thẳng cột:}\\
  Số bị trừ ở trên, số trừ ở dưới sao cho hàng đơn vị thẳng hàng đơn vị, hàng chục thẳng hàng chục.
  \item \textbf{Bước 2: Trừ từ phải sang trái:}
  \begin{itemize}\setlength{\itemsep}{0.5pt}\setlength{\leftskip}{-2.5mm}
    \item \textbf{Hàng đơn vị:} Nếu chữ số trên bé hơn chữ số dưới, mượn $1$ chục thành mười mấy để trừ, viết chữ số đơn vị, \textbf{nhớ 1 vào hàng chục của số trừ}.
    \item \textbf{Hàng chục:} Thêm $1$ vào hàng chục của số trừ, rồi lấy hàng chục số bị trừ trừ đi.
  \end{itemize}
\end{itemize}
}
\end{tcolorbox}
\end{minipage}
\hfill
\begin{minipage}[t]{0.49\textwidth}\vspace{0pt}
\begin{tcolorbox}[
  enhanced,
  colback=teal!4!white,
  colframe=teal!75!black,
  arc=2mm, boxrule=0.8pt,
  left=2mm, right=2mm, top=1.8mm, bottom=1.8mm,
  title={\footnotesize\bfseries\color{white}\faSliders\ 2. Phân tích chi tiết ví dụ $64 - 28$}
]
{\footnotesize
\begin{center}
\begin{tabular}{r@{\,}ll}
  $6$ & $4$ & \quad \textbf{Thao tác tính:} \\
$-$ & & \quad $\bullet$ $4$ không trừ được $8$, \\
  $2$ & $8$ & \quad \ \ lấy $14 - 8 = 6$, viết $\mathbf{6}$, \textbf{nhớ 1}. \\ \cline{1-2}
  $\mathbf{3}$ & $\mathbf{6}$ & \quad $\bullet$ $2$ \textbf{thêm 1 bằng 3}; \\
      &     & \quad \ \ $6 - 3 = 3$, viết $\mathbf{3}$.
\end{tabular}
\end{center}
\vspace{1mm}
$\implies$ Kết quả: $64 - 28 = \mathbf{36}$.\\
\textbf{Kiểm tra lại:}\\
$36 + 28 = 64$ (chính xác 100\%).
}
\end{tcolorbox}
\end{minipage}

\vspace{-1.5mm}
\begin{tcolorbox}[
  enhanced,
  colback=orange!4!white,
  colframe=orange!80!black,
  arc=2mm, boxrule=0.8pt,
  left=2mm, right=2mm, top=1.5mm, bottom=1.5mm,
  title={\footnotesize\bfseries\color{white}\faSitemap\ 3. Sơ đồ đối chiếu quy trình trừ có nhớ}
]
\centering
\begin{tikzpicture}[x=1.3cm, y=0.8cm]
  \node[rectangle, rounded corners=2pt, draw=violet!80!black, fill=violet!15, font=\scriptsize\bfseries] (dvi) at (0,0) {HÀNG ĐƠN VỊ\\$4 < 8 \implies 14 - 8 = \mathbf{6}$\\Viết $\mathbf{6}$, nhớ 1 vào số trừ};
  \draw[->, very thick, red!85!black] (dvi) -- node[above, font=\tiny\bfseries] {Nhớ 1 vào 2} (3.2,0);
  \node[rectangle, rounded corners=2pt, draw=teal!80!black, fill=teal!15, font=\scriptsize\bfseries] (chuc) at (4.8,0) {HÀNG CHỤC\\$2$ thêm $1 = 3$\\$6 - 3 = \mathbf{3}$};
  \draw[->, very thick, green!70!black] (chuc) -- node[above, font=\tiny\bfseries] {Ghép lại} (6.6,0);
  \node[rectangle, rounded corners=2pt, draw=green!80!black, fill=green!20, font=\normalsize\bfseries] at (7.4,0) {$\mathbf{36}$};
\end{tikzpicture}
\end{tcolorbox}

\vspace{-1.5mm}
\noindent
\begin{minipage}[t]{0.49\textwidth}\vspace{0pt}
\begin{tcolorbox}[
  enhanced,
  colback=green!4!white,
  colframe=green!60!black,
  arc=2mm, boxrule=0.8pt,
  left=2mm, right=2mm, top=1.8mm, bottom=1.8mm,
  title={\footnotesize\bfseries\color{white}\faCheckCircle\ 4. Ví dụ thực hành chuẩn mực}
]
{\footnotesize
\textbf{Đặt tính rồi tính:}\\
\begin{tabular}{r@{\,}lr@{\,}lr@{\,}l}
  $7$ & $2$ & $8$ & $5$ & $9$ & $0$ \\
$-$ & & $-$ & & $-$ & \\
  $3$ & $6$ & $4$ & $8$ & $4$ & $3$ \\ \hline
  $\mathbf{3}$ & $\mathbf{6}$ & $\mathbf{3}$ & $\mathbf{7}$ & $\mathbf{4}$ & $\mathbf{7}$
\end{tabular}
\quad
\begin{tabular}{l}
  $\bullet$ $12 - 6 = 6$; $7 - (3+1) = 3$.\\
  $\bullet$ $15 - 8 = 7$; $8 - (4+1) = 3$.\\
  $\bullet$ $10 - 3 = 7$; $9 - (4+1) = 4$.
\end{tabular}
}
\end{tcolorbox}
\end{minipage}
\hfill
\begin{minipage}[t]{0.49\textwidth}\vspace{0pt}
\begin{tcolorbox}[
  enhanced,
  colback=red!4!white,
  colframe=red!75!black,
  arc=2mm, boxrule=0.8pt,
  left=2mm, right=2mm, top=1.8mm, bottom=1.8mm,
  title={\footnotesize\bfseries\color{white}\faExclamationTriangle\ 5. Mẹo tránh bẫy trắc nghiệm}
]
{\footnotesize
\begin{itemize}\setlength{\itemsep}{1pt}\setlength{\leftskip}{-2.5mm}
  \item \textbf{Lỗi quên thêm 1 nhớ vào số trừ:}\\
  $64 - 28 = 46$ (sai vì chỉ lấy $6 - 2 = 4$). Nhớ phải lấy $6 - 3 = 3$.
  \item \textbf{Quy tắc kiểm tra thần tốc:}\\
  Sau khi trừ xong, lấy kết quả cộng lại với số trừ. Nếu đúng bằng số bị trừ thì chắc chắn $100\%$ đúng!
\end{itemize}
}
\end{tcolorbox}
\end{minipage}


\end{document}
